% ============================ WORKSHEET 2 ==================================
\worksheetheading{2}{Quantum model of the atom}
\data{$c = \SI{3.00e8}{\metre\per\second}$\quad $h = \SI{6.62e-34}{\joule\second}$\quad
$1\ \mathrm{eV} = \SI{1.6e-19}{\joule}$\quad $N_A = \SI{6.022e23}{\per\mole}$\quad
$R_H = \SI{3.29e15}{\hertz}$\quad hydrogen: $E_n = -13.6/n^2$~eV; hydrogen-like ion:
$E_n = -13.6\,Z^2/n^2$~eV. Calculator allowed.}

\wq{1}{Ordering the spectrum}
Put these radiations in order of increasing frequency: visible light, X-rays, radio waves, infrared,
$\gamma$ rays, ultraviolet, microwaves. Which one has the longest wavelength? Which photons carry the
most energy?
\begin{solution}
Radio waves $<$ microwaves $<$ infrared $<$ visible $<$ ultraviolet $<$ X-rays $<$ $\gamma$ rays.
Radio waves have the longest wavelength ($\lambda = c/\nu$); $\gamma$ photons carry the most energy
($E = h\nu$).
\end{solution}

\wq{1}{A red laser}
A laser emits red light at $\lambda = 650$~nm. Compute its frequency $\nu$ and its period $T$.
\begin{solution}
$\nu = c/\lambda = 3.00\times10^8/650\times10^{-9} = 4.62\times10^{14}$~Hz;\quad
$T = 1/\nu = \lambda/c = 2.17\times10^{-15}$~s.
\end{solution}

\wq{1}{Energy of one photon}
Compute the energy of one photon of the laser of Q2, in joules and in electronvolts.
\begin{solution}
$E = hc/\lambda = 6.62\times10^{-34}\times3.00\times10^8/650\times10^{-9} = 3.06\times10^{-19}$~J;\quad
$E = 3.06\times10^{-19}/1.6\times10^{-19} = 1.91$~eV.
\end{solution}

\wq{1}{Vocabulary of the Bohr model}
Define ground state, excited state and ionised state for the hydrogen atom, and give the energy
of the electron in each case (for the excited state, give $E_2$ and $E_3$).
\begin{solution}
Ground state: $n = 1$, lowest energy, $E_1 = -13.6$~eV. Excited state: any $n > 1$, higher energy,
e.g.\ $E_2 = -13.6/4 = -3.40$~eV, $E_3 = -13.6/9 = -1.51$~eV. Ionised state: $n\to\infty$,
$E_\infty = 0$, the electron is no longer bound to the proton.
\end{solution}

\wq{2}{Emission versus absorption}
Describe the emission spectrum and the absorption spectrum of hydrogen gas in the visible. Explain
why the lines appear at the same wavelengths in both.
\begin{solution}
Emission (gas excited by a discharge): four coloured lines on a black background at 656, 486, 434
and 410~nm. Absorption (white light through the gas): a continuous rainbow with four black lines at
the same wavelengths. Both involve the same pairs of levels: in emission an electron falls from $E_p$
to $E_n$ and emits a photon $h\nu = E_p - E_n$; in absorption a photon of the same energy lifts it from
$E_n$ to $E_p$. Same energy, same wavelength.
\end{solution}

\wq{2}{Why is permanganate purple?}
An aqueous solution of $\ce{KMnO4}$ absorbs light between about 515 and 565~nm.
(a) Which colour is absorbed? What colour does the solution appear? (b) Why is the absorption a
broad band and not thin lines as for hydrogen gas?
\begin{solution}
(a) 515--565~nm is green. The eye receives the rest of the white light, mainly red and blue--violet,
which looks purple/magenta: the complementary colour of green.\\
(b) The absorbing species is a molecular ion in a liquid: it has very many closely spaced energy levels
(vibrations, interactions with the solvent), so it absorbs a continuous range of energies. An
isolated atom in a gas has only a few well-separated levels, hence thin lines.
\end{solution}

\wq{2}{The Balmer lines from Rydberg's formula}
Using $\nu = R_H(1/n_1^2 - 1/n_2^2)$ with $n_1 = 2$, compute the wavelengths for $n_2 = 3, 4, 5, 6$
and compare them with the handout values.
\begin{solution}
$\lambda = c/\nu$:\quad
$n_2 = 3$: $\nu = 3.29\times10^{15}(\frac14-\frac19) = 4.57\times10^{14}$~Hz, $\lambda = 657$~nm;\quad
$n_2 = 4$: $6.17\times10^{14}$~Hz, 486~nm;\quad
$n_2 = 5$: $6.91\times10^{14}$~Hz, 434~nm;\quad
$n_2 = 6$: $7.31\times10^{14}$~Hz, 410~nm.
They agree with 656, 486, 434, 410~nm (the small gap for the first line comes from rounding $c$ and $R_H$).
\end{solution}

\wq{2}{Rydberg meets Bohr}
(a) Using Bohr's levels, compute the energy (in eV and J) of the photon emitted in the transition
$3\to2$, and its wavelength. (b) Show that Bohr's formula gives Rydberg's formula and that
$R_H = 13.6\ \mathrm{eV}/h$. Check the numerical value.
\begin{solution}
(a) $\Delta E = E_3 - E_2 = 13.6(\frac14 - \frac19) = 1.89$~eV $= 3.02\times10^{-19}$~J;
$\lambda = hc/\Delta E = 657$~nm (the red line).\\
(b) $h\nu = E_{n_2} - E_{n_1} = 13.6\ \mathrm{eV}\,(1/n_1^2 - 1/n_2^2)$, so
$\nu = (13.6\ \mathrm{eV}/h)(1/n_1^2 - 1/n_2^2)$: Rydberg's formula with
$R_H = 13.6\times1.6\times10^{-19}/6.62\times10^{-34} = 3.29\times10^{15}$~Hz. \checkmark
\end{solution}

\wq{3}{Series and their limits}
(a) Compute the wavelength of the first Lyman line ($2\to1$) and the Lyman series limit. In which
region are they? (b) Compute the Paschen series limit ($n_1 = 3$). (c) Explain why only the Balmer
series has lines in the visible.
\begin{solution}
(a) $2\to1$: $\nu = R_H\times\frac34 = 2.47\times10^{15}$~Hz, $\lambda = 122$~nm. Limit:
$\nu = R_H$, $\lambda = c/R_H = 91.2$~nm. Both UV.\\
(b) $\nu = R_H/9 = 3.66\times10^{14}$~Hz, $\lambda = 821$~nm (IR). All Paschen lines are beyond
821~nm.\\
(c) Lyman gaps are at least $E_2 - E_1 = 10.2$~eV (UV); Paschen gaps are at most 1.51~eV
($\lambda\ge821$~nm, IR); Balmer gaps go from 1.89~eV (656~nm) to 3.40~eV (365~nm), which covers
the visible range 400--700~nm.
\end{solution}

\wq{3}{Ionisation energy of hydrogen}
(a) What energy is needed to ionise a hydrogen atom in its ground state? Give it in eV, in J per atom
and in $\mathrm{kJ\,mol^{-1}}$. (b) What is the longest wavelength that can do it?
(c) Unit 5 lists ionisation energies in $\mathrm{kJ\,mol^{-1}}$ but leaves hydrogen blank: fill it in.
\begin{solution}
(a) $E_\infty - E_1 = 13.6$~eV $= 13.6\times1.6\times10^{-19} = 2.18\times10^{-18}$~J per atom;
$\times N_A$: $1.31\times10^{6}\ \mathrm{J\,mol^{-1}} = 1310\ \mathrm{kJ\,mol^{-1}}$.\\
(b) $\lambda = hc/E = 91.2$~nm (the Lyman limit, UV); shorter wavelengths ionise too.\\
(c) Hydrogen: about $1310\ \mathrm{kJ\,mol^{-1}}$ (the accepted value is 1312).
\end{solution}

\wq{3}{Counting emission lines}
Hydrogen atoms are excited to $n = 4$ and then return to the ground state by all possible paths.
(a) How many different lines can be emitted? List the transitions. (b) Sort them into series and
spectral regions. (c) Which ones can you see?
\begin{solution}
(a) Any pair of levels among $\{1,2,3,4\}$: $\binom42 = 6$ lines:
$4\to3$, $4\to2$, $4\to1$, $3\to2$, $3\to1$, $2\to1$.\\
(b) Lyman (to 1, UV): $4\to1$ (97~nm), $3\to1$ (103~nm), $2\to1$ (122~nm). Balmer (to 2):
$4\to2$ (486~nm), $3\to2$ (656~nm), both visible. Paschen (to 3): $4\to3$ (1876~nm, IR).\\
(c) Only $4\to2$ (blue-green) and $3\to2$ (red).
\end{solution}

\wq{3}{Which photons are absorbed?}
Hydrogen atoms in their ground state are hit by photons of energy (a) 10.2~eV, (b) 11.0~eV,
(c) 12.09~eV, (d) 15.0~eV. For each, say whether it is absorbed and what happens to the atom.
\begin{solution}
The possible gaps from $n = 1$ are $E_p - E_1 = 13.6(1 - 1/p^2)$: 10.20 ($p = 2$), 12.09 ($p = 3$),
12.75 ($p = 4$) \dots\ up to 13.6~eV (ionisation).\\
(a) Absorbed, $n = 1\to2$. (b) Not absorbed: 11.0~eV matches no gap and is below 13.6~eV.
(c) Absorbed, $n = 1\to3$. (d) Absorbed: $15.0 > 13.6$~eV, the atom is ionised and the electron leaves
with kinetic energy $15.0 - 13.6 = 1.4$~eV.
\end{solution}

\wq{4}{Identify the transition}
A line of the hydrogen emission spectrum is measured at $\lambda = 1876$~nm. Find $n_1$ and $n_2$.
Explain your method.
\begin{solution}
$\nu = c/\lambda = 3.00\times10^8/1876\times10^{-9} = 1.599\times10^{14}$~Hz, so
$1/n_1^2 - 1/n_2^2 = \nu/R_H = 0.0486$. The line is in the IR, and the Balmer limit is 365~nm, so
$n_1\ge3$. Try $n_1 = 3$: $1/n_2^2 = 1/9 - 0.0486 = 0.0625 = 1/16$, so $n_2 = 4$.
The line is $4\to3$, the first Paschen line. (With $n_1 = 4$, $1/n_2^2 = 0.0625 - 0.0486 = 0.0139$ gives
$n_2 = 8.5$, not an integer.)
\end{solution}

\wq{4}{The helium ion}
$\ce{He+}$ ($Z = 2$) has a single electron. (a) Give its energy levels $E_1$, $E_2$, $E_3$ and its
ionisation energy. (b) Compute the wavelength of its $3\to2$ transition; in which region is it?
(c) Show that the $\ce{He+}$ transition $6\to4$ has exactly the same wavelength as the red line of
hydrogen. Why is this so?
\begin{solution}
(a) $E_n = -13.6\times4/n^2 = -54.4/n^2$~eV: $E_1 = -54.4$, $E_2 = -13.6$, $E_3 = -6.04$~eV.
Ionisation energy 54.4~eV.\\
(b) $\Delta E = 54.4(\frac14-\frac19) = 7.56$~eV; $\lambda = hc/\Delta E = 164$~nm (UV).\\
(c) $6\to4$: $\Delta E = 54.4(\frac1{16}-\frac1{36}) = 13.6\times4\times\frac{5}{144} =
13.6\times\frac5{36} = 13.6(\frac14-\frac19)$, the same energy as H $3\to2$, so $\lambda = 656$~nm.
Because $Z^2/n^2 = (Z/n)^2$ and $4/6 = 2/3$, $2/4 = 1/2$: every $\ce{He+}$ level with even $n = 2k$ has the
same energy as the H level $k$.
\end{solution}

\wq{4}{Photons in a laser beam}
A green laser pointer emits 1.0~mW at 532~nm. (a) How many photons does it emit per second?
(b) How many moles of photons per hour? (c) Why do we not notice that light comes in packets?
\begin{solution}
(a) One photon: $E = hc/\lambda = 6.62\times10^{-34}\times3.00\times10^8/532\times10^{-9} = 3.73\times10^{-19}$~J.
Per second: $1.0\times10^{-3}/3.73\times10^{-19} = 2.7\times10^{15}$ photons.\\
(b) $2.7\times10^{15}\times3600/6.022\times10^{23} = 1.6\times10^{-5}$~mol per hour.\\
(c) Each packet is tiny ($\sim10^{-19}$~J) and there are $\sim10^{15}$ of them per second: the beam looks
continuous, just as water looks continuous although it is made of molecules.
\end{solution}
